\begin{theorem}[A projector and its complement]
    \label{thm:peps_binary_projection_family}
    \lean{Matrix.binaryProjectionFamily_complete}
    \leanok
    Let $D$ be a Hermitian idempotent on a finite-dimensional complex space.
    Then $D$ and $I-D$ are positive orthogonal projectors, their product in
    either order is zero, and their sum is $I$. This auxiliary identity
    supplies the two outcomes of the interference measurement in
    \cite[Section~6.6]{Schuch2010PEPS}, source lines 2605--2615.
\end{theorem}

\begin{proof}\leanok
    Hermitian idempotence makes each operator its own product with its
    adjoint, hence positive. Expanding the products proves the remaining
    identities.
\end{proof}

\begin{definition}[Transport of a commuting projection]
    \label{def:peps_scaled_projection_transport}
    \lean{Matrix.scaledProjectionTransport}
    \leanok
    Let $T$ map a finite-dimensional virtual space to the original physical
    space. Suppose $T^*T=cP$ and $TP=T$, where $c>0$. For a Hermitian
    idempotent $D$ commuting with $P$, put
    \begin{align}
        Q & =c^{-1}TDT^*.\notag
    \end{align}
    These are the finite-dimensional relations used in the local isometry
    argument of \cite[Section~6.6]{Schuch2010PEPS}, lines 2582--2615.
\end{definition}

\begin{theorem}[Complete physical projection transport]
    \label{thm:peps_scaled_projection_transport}
    \lean{Matrix.scaledProjectionTransport_properties,
        Matrix.exists_binaryProjectionFamily_transport}
    \leanok
    \uses{def:peps_scaled_projection_transport,
        thm:peps_binary_projection_family}
    Under the preceding hypotheses, $Q$ is a positive Hermitian idempotent
    and $QT=TD$. The pair $Q,I-Q$ is a complete orthogonal measurement on
    the entire original physical space, and
    \begin{align}
        QT & =TD,\qquad (I-Q)T=T(I-D).\notag
    \end{align}
    This statement concerns the coefficient map itself, independently of
    any boundary configuration or state decomposition.
\end{theorem}

\begin{proof}\leanok
    Use $T^*T=cP$, $DP=PD$, and $TP=T$ to obtain $QT=TD$.
    Hermiticity follows from that of $D$, and $D^2=D$ then gives $Q^2=Q$.
    Apply the preceding binary projection identity and expand $(I-Q)T$.
\end{proof}

\begin{theorem}[Return projection on literal pair coefficients]
    \label{thm:peps_literal_pair_return_projection}
    \lean{TNLean.PEPS.regularChargePairReturnProjection_mulVec_braidedCoefficient,
        TNLean.PEPS.regularChargePairReturnProjection_commute}
    \leanok
    \uses{thm:peps_regular_charge_pair_return_projection}
    Let $\chi$ be the character of a unitary irreducible representation of
    $G$, of dimension $d$. Put
    \begin{align}
        w(g,h) & =\chi(ph^{-1}g),\qquad
        w_k^x(g,h)=\chi(ph^{-1}xk^{-1}x^{-1}g),\notag \\
        D      & =|G|^{-2}|w\rangle\langle w|.\notag
    \end{align}
    Then $Dw_k^x=\chi(k^{-1})w/d$, and $D$ commutes with simultaneous
    left translation of the two group registers. The latter assertion
    holds for every function $\chi:G\to\mathbb C$, without an
    irreducibility or normalization assumption.
\end{theorem}

\begin{proof}\leanok
    Remove the common normalization from the established projector action.
    Simultaneous left translation fixes $w$, and hence its outer product.
\end{proof}
